% =====================================================================
%  CHAPTER 1F: 营养学基础 - 维生素与水
% =====================================================================
\chapter{营养学基础（六）：维生素与水}
\setcounter{chapter}{1}
\chapterintro{
掌握\keyword{维生素的共同特点}、\keyword{脂溶性维生素与水溶性维生素的特性对比}、\keyword{维生素缺乏发展的四个阶段}；掌握\keyword{维生素A（视黄醇）的概念、生理功能、缺乏与过量表现及食物来源}；掌握\keyword{维生素D的概念、生理功能、缺乏病及来源}；掌握\keyword{维生素C（抗坏血酸）的理化性质、生理功能、坏血病的表现及食物来源}；掌握\keyword{维生素B$_1$（硫胺素）的理化性质、TPP辅酶功能、脚气病的三型表现及食物来源}；掌握\keyword{维生素B$_2$（核黄素）的理化性质（对紫外光敏感）、口腔-生殖系统综合征（缺乏症）及食物来源}；掌握\keyword{烟酸（维生素PP/Niacin）的活性形式（NAD$^+$和NADP$^+$）、癞皮病（Pellagra）的三D症状}。熟悉维生素E、维生素K、维生素B$_6$、叶酸、维生素B$_{12}$的基本知识。了解泛酸、生物素、肉碱、牛磺酸、肌醇、柠檬素（生物类黄酮）及水的生理功能。
}

% ==================== 第一节 维生素概述 ====================
\section{维生素概述}

\begin{defbox}
\textbf{维生素（Vitamin）}是维持机体正常生理功能所必需的一类微量低分子有机化合物。它们既不是构成机体组织的原料，也不为机体提供能量，但在物质代谢中发挥不可或缺的调节作用。
\end{defbox}

\subsection{维生素共同特点}

\begin{keybox}
\textbf{维生素的共同特点：}
\begin{enumerate}[leftmargin=*]
    \item \textbf{以本体或前体形式存在于天然食物中}，必须经常由食物提供
    \item \textbf{需要量甚微}（mg或$\mu$g水平），但在调节机体代谢方面作用重要
    \item \textbf{不构成组织，也不提供能量}（有别于产能营养素）
    \item \textbf{多以辅酶或辅基的形式发挥功能}
    \item \textbf{有的具有几种结构相近、活性相同的化合物}
\end{enumerate}
\end{keybox}

\subsection{维生素的命名与分类}

\begin{keybox}
\textbf{维生素的命名：}
\begin{itemize}[leftmargin=*]
    \item 按发现顺序命名（如VitA、VitB、VitC、VitD等）
    \item 按化学结构命名（如视黄醇、硫胺素、核黄素等）
    \item 按生理功能命名（如抗坏血酸、抗干眼病维生素等）
\end{itemize}
\end{keybox}

\begin{keybox}
\textbf{维生素的分类*（按溶解性）：}

\textbf{脂溶性维生素：}
\begin{itemize}[leftmargin=*]
    \item \keyword{维生素A（视黄醇）}、\keyword{维生素D（钙化醇）}、\keyword{维生素E（生育酚）}、\keyword{维生素K（凝血维生素）}
    \item 吸收需要脂肪和胆汁的存在
    \item 可在体内储存（主要在肝脏和脂肪组织），缺乏症状出现缓慢
    \item 摄入过量可引起中毒（尤其VitA和VitD，VitK除外）
    \item 除VitK外，均有UL限制
\end{itemize}

\textbf{水溶性维生素：}
\begin{itemize}[leftmargin=*]
    \item \keyword{B族维生素}：B$_1$（硫胺素）、B$_2$（核黄素）、B$_3$（烟酸/维生素PP）、B$_5$（泛酸）、B$_6$（吡哆醇）、B$_7$（生物素）、B$_9$（叶酸）、B$_{12}$（钴胺素）
    \item \keyword{维生素C（抗坏血酸）}
\end{itemize}

\textbf{水溶性维生素的特点：}
\begin{itemize}[leftmargin=*]
    \item 体内储存量极少（B$_{12}$除外），多余者以原形或代谢产物经尿排出
    \item 缺乏症状出现较快（尤其B族维生素）
    \item 一般不引起中毒（过量摄入可通过尿液快速排泄）
    \item 饱和剂量以上几乎不增加组织浓度
\end{itemize}
\end{keybox}

\begin{exambox}
\textbf{考点提示：}两类维生素的比较是常考基础知识点。关键记忆点：脂溶性$\rightarrow$可储存、慢缺乏、易中毒（VitK除外）；水溶性$\rightarrow$储存少、快缺乏、不易中毒（尿液排泄）。
\end{exambox}

\subsection{维生素缺乏}

\begin{keybox}
\textbf{维生素缺乏的原因：}
\begin{itemize}[leftmargin=*]
    \item 维生素摄入量不足
    \item 吸收利用降低
    \item 维生素需要量相对增高
\end{itemize}

\textbf{缺乏分类：}
\begin{itemize}[leftmargin=*]
    \item 按原因分：\textbf{原发性缺乏}和\textbf{继发性缺乏}
    \item 按缺乏程度分：\textbf{临床缺乏}和\textbf{亚临床缺乏（边缘缺乏）}
\end{itemize}

\textbf{维生素缺乏发展四阶段（由轻到重）：}
\begin{enumerate}[leftmargin=*]
    \item \textbf{组织中维生素储量降低}（亚临床/边缘缺乏，无明显症状）
    \item \textbf{生化指标异常，生理功能降低}（生化检查可发现，如酶活性下降）
    \item \textbf{组织病理变化，出现临床症状和体征}（典型的缺乏病表现）
    \item \textbf{停止生命}（极端严重缺乏，可致死）
\end{enumerate}
\end{keybox}

% ==================== 第二节 维生素A ====================
\section{维生素A（视黄醇，Retinol）}

\subsection{概念与理化性质}

\begin{defbox}
\textbf{维生素A类}是指含有$\beta$-白芷酮环的多烯基结构、并具有视黄醇生物活性的一大类物质。
\end{defbox}

\begin{keybox}
\textbf{存在形式：}
\begin{itemize}[leftmargin=*]
    \item \textbf{已形成的维生素A：}动物体内具有视黄醇生物活性的维生素A称为已形成的维生素A，包括\keyword{视黄醇（Retinol）}、\keyword{视黄醛（Retinal）}、\keyword{视黄酸（Retinoic acid）}等物质
    \item \textbf{维生素A原（Provitamin A）：}黄、绿、红色植物含有类胡萝卜素，其中一部分在体内转变成维生素A的类胡萝卜素称为维生素A原，其中主要有$\alpha$-胡萝卜素、\imp{$\beta$-胡萝卜素}、$\gamma$-胡萝卜素和隐黄素4种，\imp{以$\beta$-胡萝卜素的活性最高}
\end{itemize}

\textbf{理化性质：}维生素A和胡萝卜素都对酸、碱和热稳定，一般烹调和罐头加工不易破坏，但\imp{易被氧化和受紫外线破坏}。当食物中含有磷脂、维生素E、维生素C和其它抗氧化剂时，视黄醇和胡萝卜素较为稳定，脂肪酸败可引起其严重破坏。

\textbf{单位：}
\begin{itemize}[leftmargin=*]
    \item 国际单位（IU）：1 IU 维生素A = 0.3 $\mu$g 全反式视黄醇
    \item 视黄醇当量（Retinol Equivalent, RE）
    \item 视黄醇活性当量（Retinol Activity Equivalent, RAE）
\end{itemize}
\end{keybox}

\subsection{生理功能}

\begin{keybox}
\textbf{维生素A的生理功能*：}
\begin{enumerate}[leftmargin=*]
    \item \textbf{维持正常视觉：}视网膜杆状细胞中的\keyword{视紫红质（Rhodopsin）}由视蛋白和11-顺视黄醛组成。光照$\rightarrow$视紫红质分解$\rightarrow$神经冲动$\rightarrow$视觉。暗处视紫红质再合成需要VitA持续供给。缺乏导致暗适应延长，严重者为\imp{夜盲症（Night Blindness）}
    \item \textbf{维持上皮的正常生长与分化：}维持皮肤和黏膜（呼吸道、消化道、泌尿道）的完整性
    \item \textbf{促进生长发育：}视黄酸通过核受体（RAR/RXR）调控基因表达
    \item \textbf{抑癌作用}
    \item \textbf{维持机体正常免疫功能：}促进免疫细胞分化和抗体生成
    \item \textbf{抗氧化作用：}$\beta$-胡萝卜素具有抗氧化功能，可淬灭单线态氧和清除自由基
\end{enumerate}
\end{keybox}

\subsection{维生素A缺乏与过量}

\begin{keybox}
\textbf{维生素A缺乏*：}

\textbf{长期缺乏的表现：}
\begin{itemize}[leftmargin=*]
    \item \textbf{夜盲症（Night Blindness）、干眼病（Xerophthalmia）}：暗适应能力下降$\rightarrow$夜盲症（最早出现的症状！）$\rightarrow$干眼病$\rightarrow$角膜软化$\rightarrow$角膜溃疡$\rightarrow$穿孔$\rightarrow$失明
    \item 身体多种上皮组织角化，保护内脏器官的作用减弱，抵抗力降低
    \item 影响生长发育和生殖能力
    \item 致癌
\end{itemize}

\textbf{维生素A营养状况鉴定：}
\begin{itemize}[leftmargin=*]
    \item 血清维生素A水平：参考指标
    \item 改进的相对剂量反应实验：诊断边缘状态
    \item 暗适应功能测定：适用大规模人群调查
    \item 血浆视黄醇结合蛋白：较好指标
    \item 眼结膜印迹细胞学法及眼部症状检查
\end{itemize}
\end{keybox}

\begin{tipbox}
\textbf{维生素A过量的危害：}
\begin{itemize}[leftmargin=*]
    \item 皮肤干燥、瘙痒、脱皮、脱发，出现皮疹
    \item 食欲下降，腹泻，肝脾肿大
    \item 易激动，肌肉无力，头痛，恶心，呕吐
    \item 易出血，凝血时间延长
    \item 破骨细胞活性增强，造成骨骼脆性增加，关节变形
    \item \textbf{致畸性：}孕早期过量VitA可导致胎儿畸形
    \item 注意：不包括$\beta$-胡萝卜素过量（仅引起皮肤黄染，停药后可消退）
\end{itemize}
\end{tipbox}

\subsection{视黄醇当量与参考摄入量}

\begin{keybox}
\textbf{视黄醇当量（Retinol Equivalent, RE）*：}

不同来源的维生素A折算公式：
\begin{center}
\imp{视黄醇当量 ($\mu$g RE) = 视黄醇 ($\mu$g) + $\beta$-胡萝卜素 ($\mu$g) $\times$ 0.167}
\end{center}

即：1 $\mu$g视黄醇 = 1 $\mu$g RE；1 $\mu$g $\beta$-胡萝卜素 = 0.167 $\mu$g RE = 1/6 $\mu$g RE

\textbf{膳食参考摄入量：}
\begin{itemize}[leftmargin=*]
    \item 成年男性 RNI：\imp{800 μg RE/d}
    \item 成年女性 RNI：\imp{700 μg RE/d}
    \item 孕妇早期：700 μg RE/d；中、晚期：770 μg RE/d
    \item 乳母：1300 μg RE/d
    \item UL（可耐受最高摄入量）：\imp{3000 μg RE/d}
\end{itemize}
\end{keybox}

\begin{exambox}
\textbf{维生素A的食物来源（常考）：}
\begin{itemize}[leftmargin=*]
    \item \textbf{动物来源：}动物肝脏（猪肝、牛肝、羊肝、鸡肝、鸭肝）、奶制品、鸡蛋、河蟹
    \item \textbf{植物来源（由$\beta$-胡萝卜素转化而来）：}胡萝卜、菠菜、冬苋菜、茴香、芥蓝、金针菜、芦笋、西兰花、芒果、紫菜
\end{itemize}
\end{exambox}

% ==================== 第三节 维生素D ====================
\section{维生素D（钙化醇，Calciferol）}

\subsection{概念与理化性质}

\begin{defbox}
\textbf{维生素D类}是指含环戊氢烯菲环结构、并具有钙化醇生物活性的一大类物质，包括\keyword{维生素D$_2$（麦角钙化醇，Ergocalciferol）}和\keyword{维生素D$_3$（胆钙化醇，Cholecalciferol）}。
\end{defbox}

\begin{keybox}
\textbf{理化性质：}维生素D性质较稳定，能抵抗酸、碱及氧化作用，对热也较稳定。所以通常烹调加工不会引起维生素D损失。但脂肪酸败可引起维生素D破坏。

\textbf{单位换算：}1 IU 维生素D$_3$ = 0.025 $\mu$g 维生素D$_3$；1 $\mu$g = 40 IU

\textbf{来源与活化*：}
\begin{itemize}[leftmargin=*]
    \item \textbf{VitD$_2$（麦角钙化醇）：}来源于植物（酵母、真菌经紫外线照射后麦角固醇转化）
    \item \textbf{VitD$_3$（胆钙化醇）：}\imp{皮肤中的7-脱氢胆固醇经紫外线（UVB，290$\sim$315nm）照射转化生成}（内源性来源，最重要）
\end{itemize}

\textbf{维生素D的活化过程：}
肝 $\xrightarrow{\text{25-羟化酶}}$ 25-(OH)D$_3$（肝内主要储存形式）$\xrightarrow{\text{肾1}\alpha\text{-羟化酶}}$ \imp{1,25-(OH)$_2$D$_3$（骨化三醇}，生物活性最强的形式）

1,25-(OH)$_2$D$_3$被称为"活性维生素D"，本质属于\keyword{类固醇激素}。
\end{keybox}

\subsection{生理功能}

\begin{keybox}
\textbf{维生素D的生理功能*：}

维持细胞内外钙浓度，调节钙磷代谢：
\begin{enumerate}[leftmargin=*]
    \item \textbf{促进小肠钙吸收：}诱导肠上皮细胞合成钙结合蛋白（CaBP）
    \item \textbf{促进肾小管对钙、磷的重吸收}
    \item \textbf{对骨细胞呈现多种作用：}调节骨骼的钙化与脱钙，维持血钙和血磷的正常水平
    \item \textbf{调节基因转录作用}
    \item \textbf{通过维生素D内分泌系统调节血钙平衡}
\end{enumerate}
\end{keybox}

\subsection{维生素D缺乏与过量}

\begin{keybox}
\textbf{维生素D缺乏*：}
\begin{itemize}[leftmargin=*]
    \item \textbf{儿童——佝偻病（Rickets）：}骨骼钙化不良，O形腿、X形腿、鸡胸、方颅、肋骨串珠、Harrison沟
    \item \textbf{成人——骨软化症（Osteomalacia）：}骨矿物化不良、骨痛、肌无力
    \item \textbf{老年人——骨质疏松症（Osteoporosis）：}骨量减少、骨折风险增高
\end{itemize}
\end{keybox}

\begin{tipbox}
\textbf{维生素D过多症：}
\begin{itemize}[leftmargin=*]
    \item 食欲下降，恶心腹泻
    \item 皮肤瘙痒
    \item 多尿，肾小管钙化，肾功能减退
    \item 心血管系统异常
\end{itemize}

\textbf{机体营养水平鉴定：}
\begin{itemize}[leftmargin=*]
    \item 血浆中1,25-(OH)$_2$-D$_3$的浓度
    \item 血浆中25-OH-D$_3$的浓度
\end{itemize}
\end{tipbox}

\subsection{参考摄入量与食物来源}

\begin{keybox}
\textbf{维生素D的膳食参考摄入量：}
\begin{itemize}[leftmargin=*]
    \item 成人 RNI：\imp{5 μg/d（200 IU）}
    \item 婴幼儿、儿童、孕妇、乳母、老年人 RNI：\imp{10 μg/d（400 IU）}
    \item UL（可耐受最高摄入量）：50 μg/d（2000 IU）
\end{itemize}
\end{keybox}

\begin{exambox}
\textbf{维生素D的来源：}
\begin{itemize}[leftmargin=*]
    \item \textbf{食物来源：}鱼肝油、海鱼、动物肝脏；全脂乳、蛋类、黄油、奶油、奶酪
    \item \textbf{其他来源：}\imp{适度的日光照射}（最经济、最主要的方式）
\end{itemize}
\end{exambox}

% ==================== 第四节 维生素E ====================
\section{维生素E（生育酚，Tocopherol）}

\begin{defbox}
\textbf{维生素E}是一组具有抗氧化活性的生育酚（Tocopherols）和生育三烯酚的统称，其中\keyword{$\alpha$-生育酚}的生物活性最强。
\end{defbox}

\textbf{生理功能：}最重要的功能是\keyword{抗氧化作用}——保护细胞膜多不饱和脂肪酸免受自由基攻击（链阻断型抗氧化剂），防止脂质过氧化（形成脂褐素即老年斑）；维持生殖功能；抗衰老；保护红细胞膜完整性。

\textbf{特点：}维生素E是\imp{毒性最小的脂溶性维生素}，口服大量耐受良好。

\begin{keybox}
\textbf{膳食参考摄入量：}
\begin{itemize}[leftmargin=*]
    \item 成人 AI：\imp{14 mg $\alpha$-TE/d}
    \item 乳母：17 mg $\alpha$-TE/d
\end{itemize}
\end{keybox}

\textbf{食物来源：}植物油（小麦胚芽油、葵花籽油、玉米油）、坚果、种子、绿叶蔬菜。

% ==================== 第五节 维生素K ====================
\section{维生素K（凝血维生素）}

\begin{keybox}
\textbf{维生素K的生理功能：}

维生素K是肝脏合成凝血因子所必需的物质，参与\keyword{凝血因子II（凝血酶原）、VII、IX、X}的谷氨酸残基$\gamma$-羧基化修饰，使其获得与Ca$^{2+}$结合的能力从而具有凝血活性。

缺乏导致凝血障碍、出血倾向。新生儿因肠道菌群未建立且母乳中VitK含量低，易出现\keyword{新生儿出血症}，出生后常规注射维生素K预防。
\end{keybox}

\textbf{食物来源：}绿色绿叶蔬菜（菠菜、甘蓝、西蓝花）含量最丰富；肠道细菌也能合成一部分。

% ==================== 第六节 维生素B1 ====================
\section{维生素B$_1$（硫胺素，Thiamin）}

\subsection{理化性质}

\begin{keybox}
\textbf{硫胺素的理化性质：}
\begin{itemize}[leftmargin=*]
    \item 硫胺素的盐酸盐和硝酸盐形式在干燥和酸性溶液中均稳定，但加热至其熔点（249℃）即分解
    \item 在碱性环境下易被氧化失去活性
    \item \imp{硫胺素对亚硫酸盐极为敏感}，如有亚硫酸盐存在时迅速分解成嘧啶和噻唑，并丧失其活性
\end{itemize}
\end{keybox}

\subsection{生理功能}

\begin{keybox}
\textbf{维生素B$_1$的生理功能*：}

\textbf{活性形式：}\keyword{焦磷酸硫胺素（TPP）}，主要以辅酶形式参与两类酶系统的反应：

\begin{enumerate}[leftmargin=*]
    \item \textbf{参与体内碳水化合物和能量代谢：}
    \begin{itemize}[leftmargin=*]
        \item \textbf{羧化酶}——催化$\alpha$-酮酸的氧化脱羧反应，如丙酮酸$\rightarrow$乙酰辅酶A、$\alpha$-酮戊二酸$\rightarrow$琥珀酰辅酶A（糖代谢与三羧酸循环关键环节）
        \item \textbf{转酮醇酶}——转运二碳单位的作用（磷酸戊糖途径）
    \end{itemize}
    \item \textbf{对神经系统的影响：}TPP直接激活神经细胞的氯离子通道，控制神经传导的启动
    \item \textbf{对消化系统的影响：}硫胺素可抑制胆碱酯酶，促进食欲、胃肠道的正常蠕动和消化液的分泌
\end{enumerate}
\end{keybox}

\subsection{维生素B$_1$缺乏——脚气病（Beriberi）}

\begin{keybox}
\textbf{维生素B$_1$缺乏症——"多发性神经炎/脚气病"*：}

\textbf{干性脚气病（Dry Beriberi）：}
\begin{itemize}[leftmargin=*]
    \item 神经系统：烦躁健忘、精神分散、多梦，继而出现外周神经炎、肌肉麻痹、腿麻木、有蚁走感
    \item 消化系统：胃肠蠕动减弱、便秘、食欲减退、消化不良
\end{itemize}

\textbf{湿性脚气病（Wet Beriberi）：}
\begin{itemize}[leftmargin=*]
    \item 循环系统：水肿、心动过速、心力衰竭、气喘、微血管渗透性增大
    \item 消化系统：厌食、便秘
\end{itemize}

\textbf{急性恶性脚气病（Acute Malignant Beriberi）：}
胸闷、心悸、气短、内出血、血压下降
\end{keybox}

\subsection{营养状况评价与食物来源}

\begin{keybox}
\textbf{人体营养状况评价：}
\begin{itemize}[leftmargin=*]
    \item 尿硫胺素负荷试验
    \item 硫胺素和肌酐含量比值
    \item 红细胞转酮醇酶活性系数（TPP效应）
\end{itemize}
\end{keybox}

\begin{keybox}
\textbf{参考摄入量：}
\begin{itemize}[leftmargin=*]
    \item 成年男性 RNI：\imp{1.4 mg/d}
    \item 成年女性 RNI：\imp{1.3 mg/d}
\end{itemize}
\end{keybox}

\begin{exambox}
\textbf{食物来源：}
\begin{itemize}[leftmargin=*]
    \item 粗杂粮、豆类、硬果、干酵母
    \item 瘦猪肉、动物内脏、鸡蛋黄
    \item 粮食、蔬菜、牛奶
\end{itemize}
谷类加工精度越高（碾磨越白），VB$_1$损失越严重。硫胺素集中在谷类外层和胚芽。
\end{exambox}

% ==================== 第七节 维生素B2 ====================
\section{维生素B$_2$（核黄素，Riboflavin）}

\subsection{理化性质}

\begin{keybox}
\textbf{核黄素的理化性质：}
\begin{itemize}[leftmargin=*]
    \item 核黄素的性质比较稳定，耐酸、耐热，不易氧化，但在碱性溶液中容易被破坏
    \item \imp{游离型核黄素对紫外光高度敏感}：在酸性条件下可光解为光黄素，在碱性条件下光解为光色素而丧失生物活性
    \item 黄色晶体，水溶液呈黄绿色荧光
\end{itemize}

\textbf{活性形式：}
\begin{itemize}[leftmargin=*]
    \item \keyword{黄素单核苷酸（FMN）}
    \item \keyword{黄素腺嘌呤二核苷酸（FAD）}
\end{itemize}

FMN和FAD是体内氧化还原酶类（黄素蛋白）的辅基，参与生物氧化中的氢传递，在能量代谢中起核心作用。
\end{keybox}

\subsection{维生素B$_2$缺乏}

\begin{keybox}
\textbf{维生素B$_2$缺乏与过量*：}

\textbf{典型缺乏症——"口腔-生殖系统综合征"（Oral-Genital Syndrome）：}
\begin{itemize}[leftmargin=*]
    \item \keyword{口角炎（Angular Stomatitis）}
    \item \keyword{唇炎（Cheilitis）}
    \item \keyword{舌炎（Glossitis）}：舌头红肿、乳头萎缩、地图舌
    \item \keyword{睑缘炎、结膜炎}
    \item \keyword{脂溢性皮炎（Seborrheic Dermatitis）}：鼻唇沟、眉间、耳后
    \item \keyword{阴囊皮炎}
\end{itemize}

\textbf{过量：}\imp{目前尚未见任何毒副作用。}
\end{keybox}

\begin{tipbox}
\textbf{维生素B$_2$缺乏的膳食因素：}
\begin{itemize}[leftmargin=*]
    \item 谷类加工过精（VB$_2$集中于外层和胚芽）
    \item 牛奶光照储存（VB$_2$光解破坏，故牛奶不宜用透明瓶储存）
    \item 绿色蔬菜摄入不足
    \item 动物性食物摄入不足
\end{itemize}
维生素B$_2$缺乏常与其他B族维生素缺乏同时发生。
\end{tipbox}

\subsection{营养状况评价与食物来源}

\begin{keybox}
\textbf{机体营养状况评价：}
\begin{itemize}[leftmargin=*]
    \item 红细胞谷胱甘肽还原酶活性系数
    \item 尿核黄素排出量
    \item 尿负荷试验
\end{itemize}
\end{keybox}

\begin{keybox}
\textbf{参考摄入量：}
\begin{itemize}[leftmargin=*]
    \item 成年男性 RNI：\imp{1.4 mg/d}
    \item 成年女性 RNI：\imp{1.2 mg/d}
\end{itemize}
\end{keybox}

\begin{exambox}
\textbf{食物来源：}
\begin{itemize}[leftmargin=*]
    \item 动物内脏、瘦肉、水产品（黄鳝、蟹）
    \item 奶类、蛋类
    \item 豆类、绿色蔬菜、食用菌、藻类、硬果
\end{itemize}
\end{exambox}

% ==================== 第八节 烟酸 ====================
\section{烟酸（Niacin，维生素B$_3$，维生素PP）}

\subsection{吸收与代谢}

\begin{keybox}
\textbf{烟酸的活性形式*：}
\begin{itemize}[leftmargin=*]
    \item \keyword{辅酶I——烟酰胺腺嘌呤二核苷酸（NAD$^+$）}
    \item \keyword{辅酶II——烟酰胺腺嘌呤二核苷酸磷酸（NADP$^+$）}
\end{itemize}

NAD$^+$和NADP$^+$是体内多种脱氢酶的辅酶，参与氧化还原反应和能量代谢（糖酵解、三羧酸循环、氧化磷酸化）。

\textbf{体内转化：}
色氨酸可在体内合成烟酸——\imp{约60mg色氨酸转化为1mg烟酸}。

\textbf{烟酸当量（Niacin Equivalent, NE）：}
\[
\text{NE (mg)} = \text{烟酸 (mg)} + \frac{\text{色氨酸 (mg)}}{60}
\]
\end{keybox}

\subsection{烟酸缺乏——癞皮病（Pellagra）}

\begin{keybox}
\textbf{烟酸缺乏症——癞皮病（Pellagra）*：}

烟酸缺乏症又称癞皮病（pellagra），主要损害\imp{皮肤，口、舌、胃肠道粘膜以及神经系统}。

\textbf{典型病例的"三D"症状：}
\begin{enumerate}[leftmargin=*]
    \item \keyword{Dermatitis（皮炎）：}暴露部位对称性皮炎（日照处红斑、粗糙、色素沉着）——Casal项链
    \item \keyword{Diarrhea（腹泻）：}消化不良、腹泻、舌炎
    \item \keyword{Dementia（痴呆）：}精神症状、记忆减退、焦虑、抑郁（严重时）
\end{enumerate}

\textbf{病因：}以玉米为主食且缺乏优质蛋白质的地区多发（玉米中烟酸为结合型无法吸收，且色氨酸含量低）。

\textbf{过量摄入的副作用：}皮肤发红、眼部感觉异常、高尿酸血症，偶见高血糖等。
\end{keybox}

\subsection{参考摄入量与食物来源}

\begin{keybox}
\textbf{参考摄入量：}
\begin{itemize}[leftmargin=*]
    \item 成年男性 RNI：\imp{14 mg NE/d}
    \item 成年女性 RNI：\imp{13 mg NE/d}
\end{itemize}
\end{keybox}

\begin{exambox}
\textbf{食物来源：}
\begin{itemize}[leftmargin=*]
    \item 动物肝脏、肉类、鱼类、花生、豆类
    \item 谷类中的烟酸80\%$\sim$90\%存在于麸皮中
    \item 乳类和蛋类中烟酸含量虽低但色氨酸含量丰富
\end{itemize}
\end{exambox}

% ==================== 第九节 维生素B6 ====================
\section{维生素B$_6$（吡哆醇，Pyridoxine）}

\begin{keybox}
\textbf{维生素B$_6$的生理功能与缺乏：}

\textbf{三种形式：}吡哆醇（Pyridoxine）、吡哆醛（Pyridoxal）、吡哆胺（Pyridoxamine）。

\textbf{活性形式：}\keyword{磷酸吡哆醛（PLP）}——最重要的辅酶形式。

\textbf{生理功能：}
\begin{enumerate}[leftmargin=*]
    \item 氨基酸代谢：参与转氨作用和脱羧反应（PLP是转氨酶的辅酶）
    \item 糖原代谢：参与糖原磷酸化酶作用（PLP使其活化）
    \item 参与不饱和脂肪酸代谢和神经递质（5-HT、GABA、多巴胺）的合成
\end{enumerate}

\textbf{缺乏表现：}脂溢性皮炎、舌炎、口角炎（类似于VB$_2$缺乏）；神经系统症状（易激惹、抑郁、周围神经炎）；小细胞低色素性贫血（参与血红素合成第一步$\delta$-ALA的生成）。
\end{keybox}

\begin{keybox}
\textbf{参考摄入量：}
\begin{itemize}[leftmargin=*]
    \item 成人 AI：\imp{1.2 mg/d}
\end{itemize}
\end{keybox}

\textbf{食物来源：}广泛分布于动物和植物性食物中，肉类、肝脏、全谷物、豆类、香蕉、土豆含量较高。

% ==================== 第十节 叶酸 ====================
\section{叶酸（Folate，维生素B$_9$）}

\begin{defbox}
\textbf{叶酸}属蝶酰谷氨酸，因首先在菠菜叶中分离而得名。活性形式为\keyword{四氢叶酸（Tetrahydrofolate, THF）}。
\end{defbox}

\begin{keybox}
\textbf{叶酸的生理功能与缺乏*：}

\textbf{生理功能：}
\begin{enumerate}[leftmargin=*]
    \item \textbf{"一碳单位"的载体：}THF是一碳单位（-CH$_3$、-CH$_2$、-CHO等）的载体，参与核苷酸（嘌呤、胸腺嘧啶dTMP）的生物合成
    \item 参与氨基酸代谢：同型半胱氨酸$\rightarrow$蛋氨酸的甲基化过程需要THF（N$^5$-甲基THF）作为甲基供体
    \item 参与血红蛋白的合成（与VitB$_{12}$共同作用）
\end{enumerate}

\textbf{缺乏*：}
\begin{enumerate}[leftmargin=*]
    \item \textbf{\imp{巨幼红细胞性贫血（Megaloblastic Anemia）}：}因DNA合成受阻，细胞分裂减慢，细胞体积增大但无法正常分裂。为叶酸缺乏最典型的表现
    \item \textbf{\imp{胎儿神经管畸形（Neural Tube Defects, NTD）}：}孕早期叶酸缺乏导致。包括脊柱裂（Spina Bifida）和无脑儿（Anencephaly）。\imp{围孕期（孕前至孕早期）补充400$\mu$g/d叶酸可有效预防NTD}
    \item \textbf{高同型半胱氨酸血症：}同型半胱氨酸$\rightarrow$蛋氨酸的甲基化受阻，血Hcy升高，为动脉粥样硬化的独立危险因素
\end{enumerate}
\end{keybox}

\begin{keybox}
\textbf{膳食叶酸当量（Dietary Folate Equivalent, DFE）*：}

食物中的叶酸（天然）与叶酸补充剂（合成）的生物利用率不同：
\begin{center}
\imp{DFE ($\mu$g) = 天然食物中叶酸 ($\mu$g) + 1.7 $\times$ 补充剂/强化食品中叶酸 ($\mu$g)}
\end{center}

\textbf{参考摄入量：}
\begin{itemize}[leftmargin=*]
    \item 成人 RNI：\imp{400 μg DFE/d}
    \item 孕妇：600 μg DFE/d；乳母：550 μg DFE/d
    \item UL（可耐受最高摄入量）：\imp{1000 μg DFE/d}
\end{itemize}
\end{keybox}

\textbf{食物来源：}深色绿叶蔬菜（菠菜、芥蓝）、动物肝脏、豆类、坚果、酵母。食物中叶酸易在加热和储存过程中损失。

% ==================== 第十一节 维生素B12 ====================
\section{维生素B$_{12}$（钴胺素，Cobalamin）}

\begin{keybox}
\textbf{维生素B$_{12}$的生理功能与缺乏：}

\textbf{特征：}
\begin{itemize}[leftmargin=*]
    \item 含有钴（Co）元素的维生素——唯一含金属元素的维生素
    \item \imp{仅存在于动物性食品中！}植物性食物不含维生素B$_{12}$（\keyword{严格素食者（Vegan）有缺乏风险}）
    \item 吸收需要\keyword{内因子（Intrinsic Factor, IF）}——胃壁细胞分泌
\end{itemize}

\textbf{生理功能：}
\begin{enumerate}[leftmargin=*]
    \item 作为蛋氨酸合成酶的辅酶（与叶酸共同参与同型半胱氨酸$\rightarrow$蛋氨酸）
    \item 参与神经系统的维护与髓鞘的形成
\end{enumerate}

\textbf{缺乏表现：}
\begin{enumerate}[leftmargin=*]
    \item \imp{巨幼红细胞性贫血：}与叶酸缺乏的贫血难以区分
    \item \textbf{神经系统损害：}脊髓后索和侧索退行性变（亚急性联合变性），而叶酸缺乏不引起神经系统损害——\imp{是与叶酸缺乏贫血的鉴别要点}！
\end{enumerate}
\end{keybox}

\begin{keybox}
\textbf{参考摄入量：}
\begin{itemize}[leftmargin=*]
    \item 成人 AI：\imp{2.4 μg/d}
    \item 孕妇：2.6 μg/d；乳母：2.8 μg/d
\end{itemize}
\end{keybox}

\textbf{食物来源：}动物肝脏、肉类、鱼类、蛋类、奶类。植物性食物基本不含VitB$_{12}$。

% ==================== 第十二节 维生素C ====================
\section{维生素C（抗坏血酸，Ascorbic Acid）}

\subsection{命名与理化性质}

\begin{keybox}
\textbf{维生素C的命名与理化性质：}
\begin{itemize}[leftmargin=*]
    \item 维生素C又名\keyword{抗坏血酸（Ascorbic Acid）}，为含6碳的$\alpha$-酮基内酯的弱酸，有酸味
    \item 是一种很强的还原剂
    \item \imp{与其它维生素相比，维生素C最不稳定，在有氧或碱性环境中极易氧化}
    \item L-型具有生物活性
\end{itemize}
\end{keybox}

\subsection{生理功能}

\begin{keybox}
\textbf{维生素C的生理功能*：}
\begin{enumerate}[leftmargin=*]
    \item \textbf{抗氧化作用：}保护细胞免受自由基的氧化损伤，可还原VitE的氧化形式使其重新具有活性
    \item \textbf{作为羟化过程底物和酶的辅因子：}
    \begin{itemize}[leftmargin=*]
        \item 维持\keyword{胶原蛋白的正常功能}：维持脯氨酸羟化酶和赖氨酸羟化酶的活性（羟化反应以VitC为辅因子），促进前胶原中脯氨酸和赖氨酸的羟化——是构成结缔组织、血管壁、骨骼基质的基础
        \item 参与胆固醇的羟化：促进胆固醇向胆汁酸的转化
    \end{itemize}
    \item \textbf{在铁的吸收、转运和贮备、叶酸转变为四氢叶酸等方面发挥重要作用：}
    \begin{itemize}[leftmargin=*]
        \item 将Fe$^{3+}$还原为Fe$^{2+}$，并与之形成可溶螯合物促进铁的吸收
        \item 促进叶酸转化为四氢叶酸（THF），维持叶酸的活性形式
    \end{itemize}
    \item \textbf{增强免疫功能：}促进抗体生成、增强白细胞吞噬功能
\end{enumerate}
\end{keybox}

\subsection{维生素C缺乏——坏血病（Scurvy）}

\begin{keybox}
\textbf{维生素C缺乏症——坏血病（Scurvy）*：}

因胶原蛋白合成障碍，导致结缔组织和血管壁损伤。

\textbf{早期症状：}
\begin{itemize}[leftmargin=*]
    \item 牙龈疼痛出血、鼻出血
    \item 无力、食欲减退
    \item 皮肤干燥
    \item 伤口愈合不良
    \item 血管脆性增加
\end{itemize}

\textbf{严重缺乏：}体内不断出血，导致死亡

维生素C缺乏潜伏期约3$\sim$4个月（因体内有少量储存）。
\end{keybox}

\subsection{过量、营养状况评价与食物来源}

\begin{tipbox}
\textbf{维生素C过量表现：}
\begin{itemize}[leftmargin=*]
    \item 维生素C毒性很低，但一次口服2$\sim$8g时可能会出现腹泻、腹胀
    \item 患有结石的病人，摄入量过多时可能增加尿中草酸盐的排泄，增加尿路结石的危险
\end{itemize}

\textbf{机体营养状态评价：}
\begin{itemize}[leftmargin=*]
    \item 负荷试验
    \item 血浆维生素C含量
    \item 白细胞中抗坏血酸浓度
\end{itemize}
\end{tipbox}

\begin{keybox}
\textbf{维生素C的膳食参考摄入量：}
\begin{itemize}[leftmargin=*]
    \item 成人 RNI：\imp{100 mg/d}
    \item 孕妇早期：100 mg/d；中、晚期：115 mg/d
    \item 乳母：150 mg/d
    \item UL（可耐受最高摄入量）：1000 mg/d
\end{itemize}
\end{keybox}

\begin{exambox}
\textbf{食物来源：}
\begin{itemize}[leftmargin=*]
    \item \textbf{水果类：}苜蓿、刺梨、沙棘、柚子、枣、红果、猕猴桃、酸枣、柑桔、草莓、橙子等
    \item \textbf{蔬菜类：}新鲜蔬菜（绿色和红、黄色的辣椒、菠菜、西红柿、韭菜等）
\end{itemize}
注意：维生素C极易热破坏和氧化损失，因此\imp{新鲜生食为最佳摄取方式}。
\end{exambox}

% ==================== 第十三节 其他维生素（简略） ====================
\section{其他维生素}

\subsection{泛酸（Pantothenic Acid，维生素B$_5$）}
是\keyword{辅酶A（CoA）}和酰基载体蛋白（ACP）的组成成分，在糖、脂肪和蛋白质代谢中起核心作用。广泛存在于各类食物中（Pan-thenic意为"到处存在"），缺乏罕见。成人AI：5.0 mg/d。

\subsection{生物素（Biotin，维生素B$_7$）}
是多种羧化酶（丙酮酸羧化酶、乙酰辅酶A羧化酶等）的辅酶，参与CO$_2$固定反应。缺乏极为罕见，除非长期摄入大量生鸡蛋清（含\keyword{抗生物素蛋白（Avidin）}，可与生物素紧密结合无法吸收，导致"蛋白伤害"）。成人AI：40 μg/d。

% ==================== 第十四节 类维生素 ====================
\section{类维生素}

\subsection{肉碱（Carnitine）}
在肝和肾中由赖氨酸和蛋氨酸合成。核心功能是作为长链脂肪酸进入线粒体基质进行$\beta$-氧化的载体（"脂肪酸运输穿梭系统"）。成人AI：200 mg/d。食物来源：红肉最丰富。

\subsection{牛磺酸（Taurine）}
含硫$\beta$-氨基酸，游离存在于组织中。功能：结合胆汁酸、视网膜发育、渗透压调节、抗氧化。对早产儿是条件必需营养素。

\subsection{肌醇（Inositol）}
参与细胞信号转导（肌醇磷脂途径），促进脂肪代谢。成人AI：1000 mg/d。来源：水果、豆类、谷类。

\subsection{柠檬素（Citrin / Bioflavonoids，维生素P）}
是一类具有抗氧化活性的黄酮类化合物，常与维生素C共存于自然界中。主要功能：增强毛细血管壁的韧性和通透性调节能力，协同维生素C增强其抗氧化作用；抑制透明质酸酶活性。缺乏可导致毛细血管脆性增加（易出现紫癜）。食物来源：柑橘类水果（柠檬、橙子、柚子）、葡萄、樱桃、青椒等富含维生素C的果蔬。目前尚无官方AI值，建议每日摄入量为150$\sim$200 mg。

% ==================== 第十五节 水 ====================
\section{水}

\subsection{水的分布}

人体总含水量约占体重的\imp{50\%$\sim$70\%}，随年龄增大含水量下降、随体脂增加含水量下降。
\begin{itemize}[leftmargin=*]
    \item 细胞内液：约占总体水量的 2/3
    \item 细胞外液：约占总体水量的 1/3（包括血液、组织间液、淋巴液等）
\end{itemize}

\subsection{水的生理功能}

\begin{enumerate}[leftmargin=*]
    \item \textbf{构成人体组织：}是人体含量最多的成分
    \item \textbf{作为溶剂和运载工具：}运送营养物质和代谢产物
    \item \textbf{参与物质代谢：}直接参与水解、水合等化学反应
    \item \textbf{调节体温：}通过排汗散热维持体温恒定
    \item \textbf{润滑作用：}关节液、胸腔液、腹腔液等起缓冲保护作用
\end{enumerate}

\subsection{水的平衡}

成人每日水需要量约\imp{2500 mL}。来源：饮水约1200mL、食物水约1000mL、内生水（代谢水）约300mL。排出：尿约1500mL、皮肤蒸发约500mL、肺呼出约350mL、粪便约150mL。

% ----- 复习题 -----
\reviewcontent{
\section{复习题}

\begin{enumerate}[leftmargin=*, itemsep=8pt]
    \item \textbf{【名词解释】}维生素、脂溶性维生素、水溶性维生素、视黄醇当量（RE）、$\beta$-胡萝卜素、夜盲症、佝偻病、骨化三醇1,25-(OH)$_2$D$_3$、脚气病、TPP、口腔-生殖系统综合征、烟酸当量（NE）、癞皮病（Pellagra）三D症状、膳食叶酸当量（DFE）、巨幼红细胞性贫血、坏血病（Scurvy）
    \item \textbf{【简答题】}简述维生素的共同特点。
    \item \textbf{【简答题】}比较脂溶性维生素和水溶性维生素的特点（至少4条）。
    \item \textbf{【简答题】}简述维生素缺乏发展的四个阶段。
    \item \textbf{【简答题】}简述维生素A的生理功能、缺乏的表现（由轻到重）及主要食物来源。
    \item \textbf{【简答题】}写出视黄醇当量（RE）的计算公式。
    \item \textbf{【简答题】}简述维生素D的来源、活化过程及生理功能。
    \item \textbf{【简答题】}简述维生素B$_1$的理化性质、TPP参与的酶系统和脚气病的三型表现。
    \item \textbf{【简答题】}简述维生素B$_2$的理化性质（重点：紫外光敏感性）及口腔-生殖系统综合征的临床表现。
    \item \textbf{【简答题】}简述烟酸（维生素B$_3$/PP）的活性形式及癞皮病（Pellagra）的典型"三D"症状。
    \item \textbf{【简答题】}写出烟酸当量（NE）和膳食叶酸当量（DFE）的计算公式。
    \item \textbf{【简答题】}简述维生素C的理化性质、生理功能及缺乏病（坏血病）的典型表现。
    \item \textbf{【简答题】}比较叶酸缺乏与维生素B$_{12}$缺乏引起的贫血有何异同？
\end{enumerate}

\subsection*{参考答案要点}

\begin{enumerate}[leftmargin=*, itemsep=5pt]
    \item 见本章各概念定义框。
    \item ①以本体或前体形式存在于天然食物中，必须经常由食物提供；②需要量甚微（mg或$\mu$g水平），但在调节机体代谢方面作用重要；③不构成组织，也不提供能量；④多以辅酶或辅基的形式发挥功能；⑤有的具有几种结构相近、活性相同的化合物。
    \item 脂溶性：需脂肪和胆汁协助吸收、可体内储存（肝和脂肪组织）、缺乏症状出现慢、摄入过量可引起中毒（VitK除外）、除VitK外均有UL限制。水溶性：直接吸收、体内储存极少（B$_{12}$除外）、缺乏症状出现快、一般不引起中毒（多余从尿排泄）、饱和剂量以上几乎不增加组织浓度。
    \item ①组织中维生素储量降低（亚临床/边缘缺乏）；②生化指标异常，生理功能降低；③组织病理变化，出现临床症状和体征；④停止生命（极端严重缺乏可致死）。
    \item 生理功能：①维持正常视觉（视紫红质）；②维持上皮的正常生长与分化；③促进生长发育；④抑癌作用；⑤维持机体正常免疫功能。缺乏（由轻到重）：暗适应能力下降$\rightarrow$夜盲症$\rightarrow$干眼病/上皮角化$\rightarrow$角膜软化溃疡$\rightarrow$失明。\textbf{食物来源：}动物肝脏、奶制品、鸡蛋、河蟹（已形成的VitA）；胡萝卜、菠菜、西兰花、芒果等（$\beta$-胡萝卜素）。
    \item 视黄醇当量（$\mu$g RE）= 视黄醇（$\mu$g）+ $\beta$-胡萝卜素（$\mu$g）$\times$ 0.167。即1 $\mu$g $\beta$-胡萝卜素 = 0.167 $\mu$g RE = 1/6 $\mu$g RE。
    \item 来源：VitD$_2$（植物，麦角固醇+UV）、VitD$_3$（皮肤，7-脱氢胆固醇+UVB 290$\sim$315nm）。活化：肝$\rightarrow$25-(OH)D$_3$$\rightarrow$肾$\rightarrow$1,25-(OH)$_2$D$_3$（活性形式，类固醇激素）。生理功能：促进小肠钙吸收、促进肾小管对钙磷重吸收、对骨细胞呈现多种作用（调节骨钙化与脱钙）、调节基因转录作用、维持血钙平衡。
    \item \textbf{理化性质：}盐酸盐和硝酸盐在干燥和酸性溶液中稳定，熔点249℃分解；碱性环境中易氧化失活；对亚硫酸盐极为敏感。\textbf{TPP参与：}羧化酶（$\alpha$-酮酸氧化脱羧）和转酮醇酶（二碳单位转运）。\textbf{脚气病三型：}干性（周围神经炎、感觉异常）；湿性（心血管系统损害、水肿、心衰）；急性恶性（胸闷、心悸、气短、内出血）。
    \item \textbf{理化性质：}比较稳定，耐酸耐热，不易氧化；碱性溶液中易破坏；游离型核黄素对紫外光高度敏感（酸性$\rightarrow$光黄素，碱性$\rightarrow$光色素，均丧失生物活性）。\textbf{活性形式：}FMN、FAD。\textbf{口腔-生殖系统综合征：}口角炎、唇炎、舌炎、睑缘炎、结膜炎、脂溢性皮炎、阴囊皮炎。
    \item 活性形式：辅酶I（NAD$^+$）和辅酶II（NADP$^+$），为多种脱氢酶的辅酶。体内来源：色氨酸可在体内转化为烟酸（约60mg色氨酸$\rightarrow$1mg烟酸）。癞皮病三D：Dermatitis（皮炎）——暴露部位对称性皮炎、腹泻（Diarrhea）、痴呆（Dementia）。
    \item \textbf{NE (mg)} = 烟酸 (mg) + 色氨酸 (mg)/60。\textbf{DFE ($\mu$g)} = 天然食物中叶酸 ($\mu$g) + 1.7 $\times$ 补充剂/强化食品中叶酸 ($\mu$g)。
    \item \textbf{理化性质：}又名抗坏血酸，含6碳$\alpha$-酮基内酯弱酸，强还原剂；最不稳定，有氧或碱性环境极易氧化。\textbf{生理功能：}抗氧化、作为羟化过程底物和酶的辅因子（维持胶原蛋白正常功能、参与胆固醇羟化）、促进铁吸收和叶酸活化、增强免疫功能。\textbf{坏血病：}牙龈疼痛出血、鼻出血、皮下瘀斑、伤口愈合不良、血管脆性增加；严重$\rightarrow$体内不断出血致死。
    \item \textbf{同：}均可引起巨幼红细胞性贫血。\textbf{异：}VitB$_{12}$缺乏伴有神经系统损害（脊髓后索、侧索退行性变/亚急性联合变性），而叶酸缺乏不引起神经病变；VitB$_{12}$仅存在于动物性食品中，需要内因子（IF）协助吸收。
\end{enumerate}

\newpage
}
